\section{Classical Mechanics as the Base of Generalization}
\label{cm:ba}

The textbook picture of generalization is familiar. Quantization is presented as
a recipe applied to a classical Hamiltonian system: promote canonical coordinates
$(q,p)$ to operators $(\hat q,\hat p)$, and promote the Poisson bracket to the
commutator,
\begin{equation}
  \{f,g\} \;\longmapsto\; \frac{1}{i\hbar}\,[\hat f,\hat g] \, .
\end{equation}
The inverse limit, $\hbar \to 0$, recovers classical mechanics via the WKB
approximation and Ehrenfest's theorem. Similarly, special relativity is presented
as reducing to Newtonian mechanics for $v \ll c$, and general relativity as
reducing to Newtonian gravitation in the weak-field, slow-motion limit.

Both halves of this picture presuppose that classical mechanics has already
delivered, for the system in question, a specific and unambiguous piece of
structure: a phase space with a well-defined symplectic (or Poisson) structure,
a Hamiltonian generating the correct equations of motion, and a family of
observables closed under the bracket. Only given that structure does
``quantize it'' or ``take the classical limit'' pick out a definite theory.

For a mass point moving under a potential, or for a system subject only to
holonomic constraints, this presupposition is satisfied and the textbook picture
is accurate. But it is easy to overstate how general this good case is. As soon
as one leaves it --- constraints that cannot be integrated to relations among
coordinates alone, systems that exchange energy or momentum with an unmodeled
environment, systems whose configuration space itself is not fixed once and for
all --- the classical structure needed as input to generalization becomes
ambiguous, multiple, or simply unavailable in closed form. In those cases the
question ``what is the classical mechanics of this system'', asked with no
reference to $\hbar$ or $c$ at all, already has more than one defensible answer.
That is the sense in which classical mechanics, even in its purely subordinate
role as a base, is \emph{far from trivial}: the base itself has to be constructed,
and the construction is not always unique.


\section{On the Inequivalence of the Formulations of Classical Mechanics}

The Newtonian, Lagrangian, Hamiltonian, and Hamilton--Jacobi formulations of
classical mechanics are routinely presented as equivalent packagings of a
single physical theory, related by Hamilton's principle, the Legendre
transform, and Jacobi's theorem. This claim is correct, but only under
conditions that are rarely stated explicitly and that fail for large,
physically important classes of systems. This paper collects and organizes
eight distinct ways in which the standard equivalence breaks down: the gap
between Hamilton's principle (stationarity of the action) and the popularly
stated ``principle of \emph{least} action'' (minimality), which fails beyond a
conjugate time fixed by the curvature of the potential and which can be made
arbitrarily short; the consequent failure of the boundary-value problem built
into Hamilton's principle itself to be well-posed---even for $H=p^2/2m+V(q)$, a
sufficiently steep $V$ makes the two-endpoint problem admit no solution, or
infinitely many, in contrast to the always-well-posed Newtonian initial-value
problem; the existence of singular Lagrangians, for which the Legendre
transform is not invertible and Dirac's theory of constraints is required; the
divergence, for nonholonomically constrained systems, between the
Lagrange--d'Alembert (``nonholonomic'') and variational (``vakonomic'')
equations of motion---a disagreement that is empirically decidable rather than
merely notational; the existence of Newtonian force laws that admit no
Lagrangian description at all, governed by the Helmholtz conditions; the
failure of the Legendre transform to be a global diffeomorphism even for
pointwise-regular Lagrangians, absent the stronger condition of
hyperregularity; the mechanistically identical failure of the Hamilton--Jacobi
equation's generating function to stay single-valued beyond a caustic; and an
active disagreement in the philosophy of physics literature about whether the
Lagrangian and Hamiltonian formulations posit the same geometric structure even
where their solution sets coincide exactly. We supply two fully worked
examples---a harmonic oscillator exhibiting explicit conjugate points and
boundary-value non-uniqueness, and a free particle subject to a simple
nonholonomic constraint whose nonholonomic and vakonomic equations of motion
are derived explicitly and shown to disagree already at the level of
qualitative behavior for identical initial data. We conclude with a short
taxonomy distinguishing empirical, mathematical, variational, and structural
senses of ``equivalence,'' arguing that no single one of these senses holds
uniformly across classical mechanics.

\subsection{Introduction}

Newton's \emph{Principia} (1687) presents mechanics as a theory of forces:
bodies possess positions and velocities, forces act on them, and the second law
$F = ma$ determines their subsequent motion. A century later, Lagrange's
\emph{M\'ecanique analytique} (1788) recast the subject as a variational
principle over generalized coordinates---a formulation Lagrange was famously
proud to have produced without a single diagram. Half a century after that,
Hamilton reformulated Lagrangian mechanics again, trading velocities for
momenta via what is now called the Legendre transform and turning the
second-order Euler--Lagrange equations into a first-order canonical system on
phase space; Jacobi, Poisson, and others quickly extended this into the rich
geometric theory of Hamiltonian mechanics. Hamilton went further still, showing
that the same dynamics could be extracted from a single first-order partial
differential equation for the action itself---the Hamilton--Jacobi equation---a
construction Jacobi turned into a practical method for solving Hamiltonian
systems by separation of variables.

The pedagogical line ever since has been that these are simply different
\emph{formulations} of one and the same physical theory: notational variants,
chosen for convenience, that always agree on physical predictions. Textbook
treatments typically present the passage between them---d'Alembert's principle
on one side, the Legendre transform on the other---as a purely mechanical,
always-available translation \cite{goldstein2001,arnold1989}.

This is true, but only inside a restrictive domain that is almost never made
explicit: essentially, for systems with holonomic constraints (or none),
described by a Lagrangian that is \emph{regular}, and in fact
\emph{hyperregular}, if one wants the equivalence to hold globally rather than
merely locally. Outside that domain the four formulations can fail to agree in
several logically independent ways, some of them merely technical and some of
them---most strikingly, in the case of nonholonomic constraints---genuinely
physical, in the sense that the different formulations make different,
experimentally distinguishable predictions for the same system.

This paper has two aims. First, to collect and organize these failure modes,
most of which are individually well known within their own specialist
literatures (geometric mechanics, constrained Hamiltonian dynamics, the
calculus of variations, philosophy of physics) but rarely assembled in one
place. Second, to give a completely worked, self-contained example of the
sharpest of these failures---the split between nonholonomic and vakonomic
mechanics---so that the disagreement between formulations is visible as
explicit, checkable mathematics rather than asserted by citation alone.

subsection~\ref{standard} reviews the four standard formulations---Newtonian,
Lagrangian, Hamiltonian, and Hamilton--Jacobi---and states the equivalence
theorems that license the textbook folklore, making their hypotheses explicit.
subsection~\ref{failures} works through eight distinct failure modes.
subsection~\ref{discussion} steps back and asks what ``equivalence'' between
formulations of a physical theory ought to mean in the first place, drawing on
a debate in the philosophy of physics literature that turns out to be closely
tied to the mathematics of subsection~\ref{failures}.
subsection~\ref{conclusion} concludes.

\subsection{The Four Standard Formulations}
\label{standard}

\subsubsection{Newtonian mechanics}

For a system of $N$ particles with positions $\mathbf{r}_i \in \mathbb{R}^3$,
Newton's second law reads $m_i \ddot{\mathbf{r}}_i = \mathbf{F}_i$, where
$\mathbf{F}_i$ is the total force on particle $i$. Constraints are handled by
splitting the force into an applied part and a constraint (reaction) part,
$\mathbf{F}_i = \mathbf{F}_i^{\mathrm{app}} + \mathbf{F}_i^{\mathrm{c}}$, and
postulating---this is d'Alembert's principle---that the constraint forces do no
virtual work: $\sum_i \mathbf{F}_i^{\mathrm{c}}\cdot \delta \mathbf{r}_i = 0$
for every virtual displacement $\delta\mathbf{r}_i$ compatible with the
constraints at a frozen instant of time. It is d'Alembert's principle, not
Newton's second law by itself, that supplies the bridge to the variational
formulations below; as we will see in subsection~\ref{nonholonomic}, exactly
how one operationalizes ``compatible with the constraints'' is where
nonholonomic constraints cause trouble.

\subsubsection{Lagrangian mechanics}

Let $Q$ be the configuration manifold and $TQ$ its tangent bundle, with induced
coordinates $(q^i,\dot q^i)$. A Lagrangian is a function $L : TQ \times
\mathbb{R} \to \mathbb{R}$, typically (but not necessarily) of the form $L = T
- V$, kinetic minus potential energy. Hamilton's principle asserts that
physical trajectories $q(t)$ extremize the action
\begin{equation}
S[q(\cdot)] = \int_{t_0}^{t_1} L\big(q(t),\dot q(t), t\big)\, dt
\end{equation}
among curves with fixed endpoints $q(t_0), q(t_1)$. Setting the first variation
to zero yields the Euler--Lagrange equations
\begin{equation}
\label{EL}
\frac{d}{dt}\frac{\partial L}{\partial\dot q^i} - \frac{\partial L}{\partial q^i} = 0 ,
    \qquad i = 1,\dots,n=\dim Q .
\end{equation}
For a system subject to \emph{ideal, holonomic} constraints, this is fully
equivalent to Newton's second law plus d'Alembert's principle, once $Q$ is
taken to be the constraint surface itself and the $q^i$ are independent
generalized coordinates on it.

\subsubsection{Hamiltonian mechanics}

Define the momenta $p_i = \partial L/\partial \dot q^i$. This is the
\emph{fiber derivative} $\mathbb{F}L: TQ \to T^*Q$, $(q,\dot q)\mapsto (q,p)$.
Where $\mathbb{F}L$ can be inverted for $\dot q$ in terms of $(q,p)$, one
defines the Hamiltonian by the Legendre transform,
\begin{equation}
H(q,p,t) = p_i \dot q^i - L(q,\dot q, t)\Big|_{\dot q = \dot q(q,p,t)} ,
\end{equation}
and Hamilton's canonical equations
\begin{equation}
\label{Ham}
\dot q^i = \frac{\partial H}{\partial p_i}, \qquad 
    \dot p_i = -\frac{\partial H}{\partial q^i}
\end{equation}
reproduce~\eqref{EL} on the nose. The cotangent bundle $T^*Q$ carries a
canonical symplectic form $\omega = dp_i \wedge dq^i$, giving Hamiltonian
mechanics its distinctive geometry: canonical transformations, Poisson
brackets, Liouville's theorem, and so on.

\subsubsection{The standard equivalence theorem, and its silent hypotheses}
\label{hypotheses}

The precise statement licensing the textbook claim is the following.

\begin{proposition}[Regular equivalence]
\label{prop:regular}
Let $L: TQ\times\mathbb{R}\to\mathbb{R}$ be \emph{regular}: the fiber Hessian
    $W_{ij} = \partial^2 L/\partial \dot q^i \partial \dot q^j$ is
    nondegenerate at every point of $TQ$. Then $\mathbb{F}L$ is a local
    diffeomorphism, $H$ is well defined locally, and solutions of~\eqref{EL}
    correspond, via $p = \mathbb{F}L(q,\dot q)$, bijectively and locally to
    solutions of~\eqref{Ham}. If in addition $\mathbb{F}L$ is a \emph{global}
    diffeomorphism $TQ \to T^*Q$---in which case $L$ is called
    \emph{hyperregular}---this correspondence extends to a global bijection
    between the full solution spaces of the two theories
    \cite{abrahammarsden}.
\end{proposition}

\subsubsection{Hamilton--Jacobi theory}
\label{hjtheory}

A fourth formulation, historically inseparable from Hamilton's own work,
replaces the $2n$ first-order equations~\eqref{Ham} by a single first-order
partial differential equation for a scalar function. Define \emph{Hamilton's
principal function} $S(q,t)$ as the action evaluated along the true trajectory
of the system that starts at some fixed reference event $(q_0,t_0)$ and ends at
$(q,t)$. Hamilton showed that, viewed as a function of its endpoint, $S$
satisfies the \emph{Hamilton--Jacobi equation}
\begin{equation}
\label{HJ}
\frac{\partial S}{\partial t}+H\Big(q,\frac{\partial S}{\partial q},t\Big)=0.
\end{equation}
Conversely, Jacobi showed how to run this construction backwards and recover
the canonical equations~\eqref{Ham} directly from~\eqref{HJ}, without
integrating an ordinary differential equation, provided one can find a
\emph{complete integral}: an $n$-parameter family of solutions $S(q,t;\alpha)$,
$\alpha=(\alpha_1,\dots,\alpha_n)$, of~\eqref{HJ} whose mixed Hessian
$\partial^2 S/\partial q^i\partial \alpha_j$ is everywhere nondegenerate on the
domain of interest.

\begin{proposition}[Jacobi's theorem]
\label{prop:jacobi-thm}
Let $S(q,t;\alpha)$ be a complete integral of~\eqref{HJ} on an open set
    $U\subset Q\times\mathbb{R}$, with $\det(\partial^2 S/\partial
    q^i\partial\alpha_j)\neq 0$ throughout $U$. Then the equations
\begin{equation}
\frac{\partial S}{\partial\alpha_j} = \beta_j \ (\text{constants}), 
    \qquad p_i = \frac{\partial S}{\partial q^i}
\end{equation}
implicitly define, for each choice of constants $(\alpha,\beta)$, a curve
    $(q(t),p(t))$ solving Hamilton's equations~\eqref{Ham} on $U$; every
    solution of~\eqref{Ham} passing through $U$ arises this way for a suitable
    choice of $(\alpha,\beta)$ \cite{arnold1989,goldstein2001}.
\end{proposition}

Proposition~\ref{prop:jacobi-thm} is the formal engine behind solving
Hamiltonian systems by separation of variables, and it is a genuine equivalence
result: wherever a complete integral with nondegenerate mixed Hessian exists,
the Hamilton--Jacobi and Hamiltonian pictures determine exactly the same
trajectories. The qualification ``on the domain of interest'' is doing real
work, however---subsection~\ref{caustics} shows that it cannot in general be
dropped, and that its failure is tied to the same mechanism that undermines the
``least action'' reading of Hamilton's principle in subsection~\ref{jacobi}.

Everything that follows in subsection~\ref{failures} concerns ways in which the
hypotheses of Proposition~\ref{prop:regular} or
Proposition~\ref{prop:jacobi-thm} can fail to hold, or ways in which their
being satisfied does not settle every reasonable sense in which one might want
two physical theories to be ``the same.'' Table~\ref{tab:summary} previews the
eight failure modes and where each one sits relative to these two propositions.

\begin{table}[htbp]
\centering
\small
\begin{tabular}{@{}p{2.6cm}p{4.4cm}p{5.6cm}@{}}
\textbf{Comparison} & \textbf{Condition needed for equivalence} & \textbf{What happens when it fails} \\
Hamilton's principle $\leftrightarrow$ least action & Trajectory ends before
    the first conjugate point & Action is a saddle point, not a minimum, past a
    conjugate point (\S\ref{jacobi}) \\
Newtonian (initial-value) $\leftrightarrow$ Hamilton's principle
    (boundary-value) & Travel time less than the first conjugate time &
    Boundary-value problem may have no solution or infinitely many
    (\S\ref{bvp}) \\
Newtonian $\leftrightarrow$ Lagrangian & Force law satisfies the Helmholtz
    self-adjointness conditions & No Lagrangian exists at all
    (\S\ref{helmholtz}) \\
Lagrangian (unconstrained) $\leftrightarrow$ Lagrangian under nonholonomic
    constraints & --- constraint type is itself the fork & Lagrange--d'Alembert
    and vakonomic mechanics diverge (\S\ref{nonholonomic}) \\
Lagrangian $\leftrightarrow$ Hamiltonian (local) & $L$ regular (nondegenerate
    fiber Hessian) & Dirac's constraint algorithm required (\S\ref{singular})
    \\
Lagrangian $\leftrightarrow$ Hamiltonian (global) & $L$ hyperregular
    ($\mathbb{F}L$ a global diffeomorphism) & $\mathbb{F}L$ fails to be onto or
    one-to-one globally (\S\ref{global}) \\
Hamiltonian $\leftrightarrow$ Hamilton--Jacobi & Complete integral exists with
    nondegenerate mixed Hessian & Generating function becomes multivalued at a
    caustic (\S\ref{caustics}) \\
``The same'' system, different representatives & --- & Gauge freedom: many
    Lagrangians/Hamiltonians per force law (\S\ref{nonunique}) \\
\end{tabular}
\caption{The eight failure modes discussed in subsection~\ref{failures}, and
    where each sits relative to the regular equivalence theorem
    (Proposition~\ref{prop:regular}) or Jacobi's theorem
    (Proposition~\ref{prop:jacobi-thm}).}
\label{tab:summary}
\end{table}

Figure~\ref{fig:map} maps the Newtonian--Lagrangian--Hamiltonian chain
specifically, with the corresponding points of failure marked; the three
further failure modes concerning the variational character of the action, the
well-posedness of its boundary-value problem, and the global validity of the
Hamilton--Jacobi correspondence (\S\ref{jacobi}, \S\ref{bvp}, \S\ref{caustics})
are internal to the Lagrangian and Hamiltonian formulations respectively and
are not separately diagrammed.

\begin{figure}[htbp]
\centering
\begin{tikzpicture}[
    box/.style={draw, rounded corners, minimum width=3.1cm, minimum height=1.05cm, align=center, font=\small},
    lbl/.style={font=\scriptsize, align=center, text width=4.6cm, text=red!55!black}
]
\node[box] (N) at (0,0)   {Newtonian\\ $F=ma$};
\node[box] (L) at (5.4,0) {Lagrangian\\ $(TQ,\,L)$};
\node[box] (H) at (10.8,0){Hamiltonian\\ $(T^*Q,\,\omega,\,H)$};

\draw[-{Latex[length=2.2mm]}, thick] (N) -- (L) node[midway, above, font=\small\itshape] {d'Alembert};
\draw[-{Latex[length=2.2mm]}, thick] (L) -- (H) node[midway, above, font=\small\itshape] {Legendre transform};

\node[lbl] at (2.7,-1.7) {fails: no Lagrangian exists (\S\ref{helmholtz}); nonholonomic constraints split into two inequivalent theories (\S\ref{nonholonomic})};
\node[lbl] at (8.1,-1.7) {fails: $L$ singular (\S\ref{singular}); $L$ regular but not hyperregular (\S\ref{global})};
\end{tikzpicture}
\caption{The standard chain of formulations, annotated with the conditions under which each arrow fails to deliver an equivalence.}
\label{fig:map}
\end{figure}

\subsection{Where Exact Equivalence Fails}
\label{failures}

\subsubsection{Extremal versus minimal: Jacobi's conjugate-point theorem}
\label{jacobi}

The equivalence between Hamilton's principle and the Euler--Lagrange
equations~\eqref{EL} is exact and unconditional: $\delta S=0$ for all
admissible variations if and only if~\eqref{EL} holds along the trajectory.
This is a first-order, local statement about the action, and none of
subsection~\ref{standard} depends on anything more. But the popular name for
Hamilton's principle---the ``principle of \emph{least} action''---asserts
something stronger: that the true trajectory does not merely make $S$
stationary but \emph{minimizes} it among nearby competitors with the same
endpoints. This stronger claim concerns the second variation, not the first,
and it is generally false, even for the most elementary holonomic,
unconstrained, regular Lagrangian system one could write down
\cite{gray2009,graytaylor2007}.

For a one-dimensional natural Lagrangian $L(x,\dot x) = \tfrac12 m\dot x^2 -
V(x)$, consider a perturbation $x(t)+\varepsilon\eta(t)$ of a solution $x(t)$
of~\eqref{EL}, with $\eta(t_0)=\eta(t_1)=0$. The second variation of the action
is
\begin{equation}
\delta^2 S[\eta] = \int_{t_0}^{t_1} \Big[ m\dot\eta(t)^2 - V''(x(t))\,\eta(t)^2 \Big]\, dt ,
\end{equation}
and the trajectory is a genuine local minimum of $S$ if and only if
$\delta^2S[\eta]>0$ for every admissible $\eta\not\equiv0$. Whether this holds
is governed by the \emph{Jacobi accessory equation}, obtained by
linearizing~\eqref{EL} about $x(t)$,
\begin{equation}
\label{jacobi-eq}
m\ddot\eta + V''(x(t))\,\eta = 0 ,
\end{equation}
together with $\eta(t_0)=0$. A time $t_c>t_0$ is a \emph{conjugate point} (or
\emph{kinetic focus}) to $t_0$ along the trajectory if the nontrivial solution
of~\eqref{jacobi-eq} with $\eta(t_0)=0$ vanishes again at $t_c$.

\begin{proposition}[Jacobi's necessary condition]
\label{prop:jacobi-necessary}
The trajectory $x(t)$ furnishes a local minimum of the action on $[t_0,t_1]$
    only if $t_1$ occurs before the first conjugate point to $t_0$. If $t_1$
    lies beyond the first conjugate point, $S$ is a saddle point at $x(t)$:
    nearby trial trajectories exist with strictly smaller action, even though
    $x(t)$ exactly satisfies~\eqref{EL} \cite{gelfandfomin1963,gray2009}.
\end{proposition}

\begin{example}[Harmonic oscillator]
\label{sho}
Take $V(x)=\tfrac12 kx^2$, so $V''\equiv k$ and~\eqref{jacobi-eq} becomes
    $\ddot\eta+\omega^2\eta=0$ with $\omega=\sqrt{k/m}$. The solution with
    $\eta(t_0)=0$ is $\eta(t)=\sin\!\big(\omega(t-t_0)\big)$, which next
    vanishes at $t_c=t_0+\pi/\omega$---exactly half a period. By
    Proposition~\ref{prop:jacobi-necessary}, a segment of oscillator trajectory
    genuinely minimizes the action if its duration is less than $\pi/\omega$,
    and is merely a stationary saddle point if its duration exceeds
    $\pi/\omega$.
\end{example}

Example~\ref{sho} answers, quantitatively, the question of how short is ``short
enough.'' The conjugate-point time $\pi/\omega=\pi\sqrt{m/k}$ is set entirely
by the curvature $V''=k$ of the potential and shrinks without bound as $k$
grows: no fixed time interval is safe for every system, because a sufficiently
steeply curved potential will already have produced a conjugate point within
it. More generally, conjugate points arise only where $V''>0$ along the
trajectory; for potentials with $V''(x)\leq 0$ for all $x$ (a free particle,
uniform gravity, an inverted quadratic barrier) no conjugate point ever forms,
and the action is minimized for arbitrarily long times \cite{gray2009}.
Whether ``the action is minimized'' is therefore not a property of Hamilton's
principle in general---it is a property of the specific potential, and
specifically of how sharply it curves.

It is worth being precise about what this does and does not show. Hamilton's
principle, in its correct form ($\delta S=0$), remains exactly equivalent to
Newton's second law throughout; nothing here contradicts
Proposition~\ref{prop:regular}. What fails is the stronger, popularly stated
identification of Hamilton's principle with a principle of \emph{least}
action---and it fails already in the cleanest setting available anywhere in
this paper: a single unconstrained particle, regular and hyperregular
Lagrangian, smooth potential. That is precisely why this failure mode belongs
at the head of this subsection, ahead of the more exotic constructions that
follow.

\subsubsection{Non-existence and non-uniqueness: the boundary-value problem}
\label{bvp}

subsection~\ref{jacobi} asked whether a trajectory already known to
satisfy~\eqref{EL} minimizes the action. A logically prior question is whether
such a trajectory exists and is unique at all, given the data Hamilton's
principle is actually stated in terms of: not an initial position and velocity,
but two fixed \emph{endpoints}, $q(t_0)=q_A$ and $q(t_1)=q_B$---exactly the
boundary conditions built into the action functional in
subsection~\ref{standard}. This is a fundamentally different mathematical
problem from the Newtonian initial-value problem. Given $(q_0,\dot q_0)$ at a
single instant, standard ODE theory (Picard--Lindelöf) guarantees a unique
solution of~\eqref{EL}, at least locally in time, for any reasonably smooth
force law: the ``most boring'' Hamiltonian, $H=p^2/2m+V(q)$, is always
well-posed as an initial-value problem. Fixing $q(t_0)$ and $q(t_1)$ instead
offers no such guarantee. The boundary-value problem can have no solution, a
unique solution, or infinitely many, and which of these holds is governed by
exactly the conjugate points of subsection~\ref{jacobi}.

Fix $t_0$ and $q_A=q(t_0)$, and let $x(t,v_0)$ denote the unique initial-value
solution with $\dot x(t_0)=v_0$. The \emph{shooting map} $v_0\mapsto
x(t_1,v_0)$ determines which endpoints are reachable in time $T=t_1-t_0$ and
how many initial velocities reach a given one. Its derivative $\partial
x(t,v_0)/\partial v_0$ satisfies~\eqref{EL} linearized about the
trajectory---for $H=p^2/2m+V(q)$ this is exactly the Jacobi accessory
equation~\eqref{jacobi-eq}, with initial data $\eta(t_0)=0$, $\dot\eta(t_0)=1$.
Consequently
\begin{equation}
\frac{\partial x(t_1,v_0)}{\partial v_0} = 0
    \quad\Longleftrightarrow\quad t_1 \text{ is a conjugate point to } t_0 ,
\end{equation}
so the shooting map is a local diffeomorphism precisely up to the first
conjugate time, and can fold---losing injectivity, surjectivity, or both---at
and beyond it.

\begin{example}[Harmonic oscillator boundary-value problem]
\label{sho-bvp}
Continuing Example~\ref{sho}, set $t_0=0$, $q_A=0$. The initial-value solution
    is $x(t,v_0)=(v_0/\omega)\sin(\omega t)$, so the shooting map to time $T$
    is $v_0\mapsto (v_0/\omega)\sin(\omega T)$, and
\begin{equation}
v_0 = \frac{\omega\, q_B}{\sin(\omega T)}
\end{equation}
solves $x(0)=0$, $x(T)=q_B$ whenever $\sin(\omega T)\neq0$: a unique solution,
    as expected for $T$ short of the first conjugate time $\pi/\omega$. But at
    $T=n\pi/\omega$ for a positive integer $n$---an integer multiple of the
    conjugate time---$\sin(\omega T)=0$ identically, and the shooting map
    collapses to $v_0\mapsto 0$ regardless of $v_0$: for $q_B\neq0$ \emph{no}
    solution exists, while for $q_B=0$ \emph{every} $v_0\in\mathbb{R}$ solves
    the boundary-value problem, a one-parameter family of distinct true
    trajectories connecting the same two events \cite{gray2009}.
\end{example}

Example~\ref{sho-bvp} shows non-existence and non-uniqueness sitting on
opposite sides of the same coincidence, both triggered by the same
conjugate-time condition, and both occurring for the most elementary confining
potential there is. The critical travel times $n\pi/\omega$ are evenly spaced
only because the harmonic oscillator's own equation of motion is already
linear---so it coincides with its own Jacobi equation, and conjugate points
recur exactly periodically. For a general potential, the accessory
equation~\eqref{jacobi-eq} has a time-varying coefficient $V''(x(t))/m$ and the
conjugate times need not be evenly spaced, but the qualitative picture
survives, and can be more dramatic still: \cite{gray2009} exhibits a quartic
oscillator, $V(x)=\tfrac14 Cx^4$, in which a whole family of distinct true
trajectories leaving a common starting event recross one another at a common
later event, producing multiple genuinely different---not merely differently
scaled---solutions of the same boundary-value problem.

The dependence on how fast the potential curves is exactly the one identified
in subsection~\ref{jacobi}: since $\omega=\sqrt{k/m}$ for the harmonic case,
the critical travel times $n\pi/\omega = n\pi\sqrt{m/k}$ shrink as the
curvature $k$ grows, so for \emph{any} fixed travel time $T$, a steep enough
potential well places $T$ at or past a degenerate configuration. ``The most
boring system there is,'' $H=p^2/2m+V(q)$, therefore fails to have a well-posed
boundary-value problem in exactly the regime---strongly confining, rapidly
growing $V$---that most textbook potentials are drawn from.

This multiplicity is not a curiosity confined to classical mechanics: it is
exactly what the Feynman path-integral formulation of quantum mechanics must
sum over. The semiclassical propagator between $q_A$ and $q_B$ is, in general,
a sum of contributions from \emph{every} classical trajectory solving the
boundary-value problem, each carrying a phase determined by its action together
with a further phase set by the \emph{Maslov index}---the number of conjugate
points the trajectory has crossed \cite{gray2009}. The Maslov index is, in
this sense, a bookkeeping device for exactly the multiplicity established in
this subsection: it is needed because the classical boundary-value problem it
quantizes is not, in general, uniquely solvable.

\subsubsection{Singular Lagrangians and constrained Hamiltonian dynamics}
\label{singular}

A Lagrangian is \emph{singular} (or degenerate) when its fiber Hessian $W_{ij}
= \partial^2 L/\partial \dot q^i \partial \dot q^j$ is degenerate somewhere on
$TQ$. In that case $\mathbb{F}L$ is not a local diffeomorphism there: the
momenta $p_i = \partial L/\partial \dot q^i$ are not independent functions of
the velocities, and one cannot solve for all of $\dot q$ in terms of $(q,p)$.
Concretely, degeneracy of $W_{ij}$ means the map $\mathbb{F}L$ satisfies one or
more relations
\begin{equation}
\phi_a(q,p) = 0, \qquad a = 1,\dots,k,
\end{equation}
purely among the momenta and coordinates---\emph{primary constraints}, in
Dirac's terminology---that hold identically on the image of $\mathbb{F}L$,
before any equations of motion are imposed.

The cleanest illustration is to note that Hamiltonian mechanics is itself,
formally, a maximally singular special case of Lagrangian mechanics. Treat
$(q,p)$ as $2n$ independent coordinates on an auxiliary configuration space and
consider the first-order ``Hamiltonian-form'' Lagrangian
\begin{equation}
L_1(q,p,\dot q,\dot p) = p_i \dot q^i - H(q,p).
\end{equation}
Its Euler--Lagrange equations reproduce~\eqref{Ham} exactly, so in one sense
Hamiltonian mechanics \emph{is} a special case of Lagrangian mechanics on a
bigger space. But the fiber Hessian of $L_1$ with respect to $(\dot q,\dot p)$
vanishes identically---$L_1$ is linear, not quadratic, in the velocities---so
$L_1$ is about as singular as a Lagrangian can be, and the naive
Legendre-transform machinery of subsection~\ref{hypotheses} simply does not
apply to it. Recovering Hamilton's equations from $L_1$ requires running the
full constrained-Hamiltonian algorithm below, not the regular Legendre
transform.

The general treatment of singular Lagrangians is due to Dirac and Bergmann
\cite{dirac1964}. One passes to the \emph{primary constraint surface}
$\Sigma_1 \subset T^*Q$ cut out by the $\phi_a$, and checks whether the primary
constraints are preserved under time evolution generated by a (now only
partially determined) Hamiltonian; if not, this consistency requirement
generates \emph{secondary constraints}, and the process is iterated until a
stable constraint surface is reached. Constraints are further classified as
\emph{first class} (those with vanishing Poisson brackets with all other
constraints, which typically signal gauge freedom---directions in phase space
that are not physically distinguishable) or \emph{second class} (which can be
eliminated using the Dirac bracket, an amended Poisson bracket adapted to the
constraint surface) \cite{henneauxteitelboim1992}. The upshot is that for a
singular $L$, the physical phase space is not $T^*Q$ but a further-reduced
submanifold, of dimension strictly smaller than $2n$ in the presence of
first-class constraints, and the map from Lagrangian solutions to Hamiltonian
solutions is not the clean bijection of Proposition~\ref{prop:regular} but a
considerably more elaborate correspondence, often carrying residual gauge
freedom on the Lagrangian side that has no counterpart in a naively
written-down Hamiltonian system.

Singular Lagrangians are not a pathology confined to toy examples: they are the
generic situation for gauge theories. Electromagnetism, Yang--Mills theory, and
general relativity all have singular Lagrangians, precisely because the local
gauge symmetry (arising from Noether's second theorem) forces degeneracy of the
fiber Hessian. For essentially all of fundamental field theory, then, the
``equivalence'' between the Lagrangian and Hamiltonian pictures is not the
elementary Legendre-transform story of subsection~\ref{hypotheses} at all, but
the considerably more delicate Dirac--Bergmann correspondence.

\subsubsection{Nonholonomic constraints: Lagrange--d'Alembert versus vakonomic
mechanics}
\label{nonholonomic}

The sharpest failure of equivalence---sharp enough to be settled experimentally
rather than merely by definitional fiat---occurs already at the level of the
passage from Newtonian to Lagrangian mechanics, once the constraints are
\emph{nonholonomic}: constraints on velocities that cannot be integrated to
constraints on configuration alone. The rolling of a coin or a ball without
slipping is the standard example.

For holonomic constraints there is no ambiguity in how to build a variational
principle: one simply restricts the configuration space $Q$ to the constraint
surface and applies Hamilton's principle there. For nonholonomic constraints,
given as a distribution $\mathcal{D}\subset TQ$ (a smoothly varying subspace of
admissible velocities), there turn out to be two inequivalent ways to
generalize the construction, and classical mechanics contains no a priori
argument for preferring one over the other on variational grounds alone:

\begin{itemize}
\item \textbf{Nonholonomic (Lagrange--d'Alembert) mechanics.} One directly
    generalizes d'Alembert's principle: the physical trajectory satisfies the
        constraint at every instant, and virtual displacements $\delta q$ used
        in varying the action are required to lie in $\mathcal{D}$ pointwise
        along the trajectory, \emph{without} imposing the constraint as a
        restriction on the class of curves being compared. This is the
        classical, and---as discussed below---experimentally correct,
        treatment.
\item \textbf{Vakonomic mechanics.} (The name, coined by Kozlov, abbreviates
    ``variational axiomatic kind'' of mechanics; \cite{kozlov1983,akn2006}.)
        One instead builds the constraint directly into the variational
        problem, extremizing the action over the class of curves that satisfy
        the constraint at every instant, in the ordinary sense of a constrained
        extremum problem (formally, via a Lagrange multiplier function adjoined
        to the Lagrangian).
\end{itemize}

These two constructions look superficially similar---both are ``variational''
and both use Lagrange multipliers to enforce the same constraint---but they
generically yield \emph{different} equations of motion for the same Lagrangian
and the same constraint. \cite{lewismurray1995} worked this comparison out in
detail for a ball rolling without slipping on a rotating table, deriving both
sets of equations and testing each against a physical experiment; they found
good agreement between the nonholonomic equations (once realistic friction is
included) and the observed motion, and were unable to reproduce the
observations with the vakonomic equations. The nonholonomic/vakonomic split is
accordingly not a matter of taste or convenience: it is a case in which two
ways of extending the Lagrangian formalism to a wider class of systems make
different, checkable predictions, and only one of them is the physically
correct generalization of Newtonian mechanics with ideal constraint forces
\cite{neimarkfufaev1972,bloch2003}.

\subsubsection{A fully worked example}

To make the disagreement completely explicit, consider the simplest possible
illustration: a free particle of unit mass in $Q = \mathbb{R}^3$, coordinates
$(x,y,z)$, with Lagrangian
\begin{equation}
L(q,\dot q) = \tfrac12\big(\dot x^2 + \dot y^2 + \dot z^2\big),
\end{equation}
subject to the nonholonomic constraint
\begin{equation}
\label{constraint}
\dot z = y\,\dot x .
\end{equation}
This constraint is genuinely nonholonomic: writing it as $\omega(\dot q) = 0$
with the one-form $\omega = -y\,dx + dz$, one computes $d\omega = dx\wedge dy$
and $\omega\wedge d\omega = dx\wedge dy \wedge dz \neq 0$, so the Frobenius
integrability condition fails and $\omega$ admits no integrating factor. There
is no function $f(x,y,z)$ whose level sets reproduce~\eqref{constraint}.

\paragraph{Nonholonomic (Lagrange--d'Alembert) dynamics.}
Admissible virtual displacements satisfy $\omega(\delta q)=0$, i.e.\ $\delta z
= y\,\delta x$, with $\delta x,\delta y$ free. Since $L$ has no potential,
$\partial L/\partial q^i \equiv 0$ and $\partial L/\partial \dot q^i = \dot
q^i$, so the Lagrange--d'Alembert principle $\int (\ddot x\,\delta x+\ddot
y\,\delta y+\ddot z\,\delta z)\,dt = 0$ becomes, after substituting $\delta z =
y\,\delta x$,
\begin{equation}
\int \big[(\ddot x + y\ddot z)\,\delta x + \ddot y\,\delta y\big]\,dt = 0
\end{equation}
for arbitrary independent $\delta x(t),\delta y(t)$ vanishing at the endpoints.
Hence
\begin{align}
\ddot y &= 0, \label{nh1}\\
\ddot x &= -y\,\ddot z , \label{nh2}
\end{align}
together with the constraint~\eqref{constraint} itself, imposed on the physical
trajectory. Equation~\eqref{nh1} says that \emph{no constraint force ever acts
in the $y$-direction}: $y(t) = y_0 + v_0 t$ exactly, for all time, where $v_0 =
\dot y(0)$. Differentiating the constraint and using~\eqref{nh1}--\eqref{nh2}
gives a closed equation for $\ddot z$ and hence, after eliminating $\ddot z$,
the single first-order relation
\begin{equation}
\frac{d}{dt}\ln \dot x = -\,\frac{v_0\, y(t)}{1+y(t)^2},
\end{equation}
which integrates (using $dy = v_0\,dt$) to the closed form
\begin{equation}
\label{nh-solution}
\dot x(t) = \dot x(0)\,\sqrt{\frac{1+y_0^2}{1+y(t)^2}}, 
    \qquad y(t) = y_0 + v_0 t ,
\end{equation}
with $z(t)$ then recovered from $\dot z = y\dot x$. The nonholonomic prediction
is thus fully explicit: $y$ moves in a straight line at constant velocity no
matter what $x$ and $z$ are doing, while $\dot x$ decays monotonically to zero
as $|y|\to\infty$.

\paragraph{Vakonomic dynamics.}
Now instead adjoin the constraint to the Lagrangian with a multiplier function
$\mu(t)$ treated as an additional independent variable,
\begin{equation}
L'(x,y,z,\mu,\dot x,\dot y,\dot z) = \tfrac12(\dot x^2+
    \dot y^2+\dot z^2) + \mu(\dot z - y\dot x),
\end{equation}
and extremize freely over $x,y,z,\mu$. The Euler--Lagrange equations for
$x,y,z$ are
\begin{align}
\ddot x &= \dot\mu\, y + \mu\,\dot y, \label{vak1}\\
\ddot y &= -\mu\,\dot x, \label{vak2}\\
\ddot z &= -\dot\mu, \label{vak3}
\end{align}
while the Euler--Lagrange equation for $\mu$ simply returns the
constraint~\eqref{constraint}. Differentiating the constraint and
substituting~\eqref{vak1} and~\eqref{vak3} yields a first-order ODE governing
the multiplier itself,
\begin{equation}
\dot\mu = -\,\frac{\dot y\,(\dot x + \mu y)}{1+y^2} .
\end{equation}

Two qualitative features distinguish this from the nonholonomic case, and both
are visible directly from the equations rather than from any numerical
simulation:

\begin{enumerate}
\item By~\eqref{vak2}, $\ddot y = -\mu\dot x$ is generically \emph{nonzero}:
    unless $\mu\equiv 0$ or $\dot x \equiv 0$ throughout the motion, $y(t)$ is
        not the straight line $y_0+v_0t$ predicted by~\eqref{nh-solution}. For
        identical Lagrangian, identical constraint, and identical initial data
        $(q(0),\dot q(0))$ with $\dot y(0)\neq 0$ and $\dot x(0)\ne0$, the two
        theories therefore make different predictions for $y(t)$, checkable in
        principle from the very first instant of nonzero acceleration.
\item In the nonholonomic case the multiplier $\lambda := \ddot z = \dot y \dot
    x/(1+y^2)$ is an \emph{algebraic} function of the instantaneous state
        $(q,\dot q)$---exactly what one expects of an ideal, workless reaction
        force, which by definition depends only on the instantaneous kinematic
        constraint, not on the system's history. In the vakonomic case, by
        contrast, $\mu$ obeys its own first-order differential equation and is
        therefore a genuine additional \emph{dynamical} degree of freedom,
        coupled to the rest of the system, rather than an instantaneous
        reaction force. This is a structural, not merely quantitative,
        difference between the two theories.
\end{enumerate}

\begin{remark}
The special case $v_0 = \dot y(0) = 0$ is degenerate in a way that is itself
    instructive: the nonholonomic multiplier $\lambda = \dot y\dot x/(1+y^2)$
    vanishes identically, and the two theories briefly coincide ($\ddot x=\ddot
    y=0$ in both). Cases of this kind---where the two dynamics agree on
    special, lower-dimensional families of initial data or for constraints with
    extra symmetry---are known and have been studied systematically
    \cite{akn2006}; they do not rescue the general equivalence, but they do
    explain why the discrepancy is easy to miss if one only ever checks a
    single convenient trajectory.
\end{remark}

\begin{proposition}
\label{prop:nonequiv}
For the Lagrangian and constraint above, and for generic initial data with
    $\dot x(0)\neq 0$ and $\dot y(0)\neq 0$, the nonholonomic and vakonomic
    equations of motion have distinct solutions. In particular, the
    nonholonomic solution satisfies $\ddot y(t) \equiv 0$ for all $t$, while
    the vakonomic solution satisfies $\ddot y(0) = -\mu(0)\dot x(0)$, which is
    nonzero for generic $\mu(0)$.
\end{proposition}

This is about as direct a counterexample to ``the formulations of classical
mechanics are equivalent'' as one could ask for: identical starting data fed
into two formalisms that are each individually coherent, each individually
variational in flavor, and each individually a natural-looking generalization
of Hamilton's principle to constrained systems, yield different physics.

\subsubsection{The inverse problem: not every Newtonian system is Lagrangian}
\label{helmholtz}

Even in the unconstrained case, the passage from Newtonian to Lagrangian
mechanics is not guaranteed to exist at all. Given a system of second-order
ODEs $\ddot q^i = f^i(q,\dot q,t)$---an arbitrary Newtonian force law, divided
by mass---one can ask whether it arises as the Euler--Lagrange equations of
\emph{some} Lagrangian. This is the \emph{inverse problem of the calculus of
variations}. The answer is governed by the \emph{Helmholtz conditions}: the
linearized (Fr\'echet) differential operator associated with the equations of
motion must be self-adjoint. \cite{douglas1941} gave a complete classification
of when a Lagrangian exists in the two-degree-of-freedom case, building on
earlier work of Hirsch and Mayer; the general $n$-dimensional theory remains an
active area of research in the geometric calculus of variations.

The familiar case that fails the Helmholtz conditions is linear
velocity-dependent damping. The damped harmonic oscillator,
\begin{equation}
m\ddot x + \gamma \dot x + k x = 0 ,
\end{equation}
is not the Euler--Lagrange equation of any Lagrangian of the ordinary
autonomous, single-particle form $L(x,\dot x)$: the friction term breaks the
self-adjointness that the Helmholtz conditions require. This is not merely a
failure of imagination; it reflects the fact that a Hamilton's-principle
Lagrangian for a single degree of freedom automatically conserves a symplectic
volume (Liouville's theorem), while a damped oscillator's phase-space volume
visibly contracts as the system loses energy to friction. One can restore a
variational description, but only at a structural cost: by doubling the degrees
of freedom (the Bateman construction, pairing the damped oscillator with a
time-reversed, amplified partner), by using an explicitly time-dependent
Lagrangian such as the Caldirola--Kanai form $L = e^{\gamma t/m}\big(\tfrac12
m\dot x^2 - \tfrac12 kx^2\big)$, or by admitting non-standard
(non-quadratic-in-velocity) Lagrangians. Each of these routes reproduces the
correct equation of motion, but none of them is the ``obvious,''
structure-preserving Lagrangian one would have guessed, and none is unique.

The general moral is that Newtonian mechanics, understood as ``any second-order
ODE system one likes,'' is strictly more permissive than Lagrangian mechanics:
the class of Lagrangian-representable force laws is a proper,
Helmholtz-constrained subset of the class of Newtonian force laws, unless one
is willing to pay for a Lagrangian description with auxiliary variables or
explicit time dependence.

\subsubsection{Global obstructions: regularity is not hyperregularity}
\label{global}

Proposition~\ref{prop:regular} distinguishes \emph{regularity} (the fiber
Hessian is nondegenerate at each point) from \emph{hyperregularity} (the fiber
derivative $\mathbb{F}L$ is a diffeomorphism of the whole of $TQ$ onto the
whole of $T^*Q$). The distinction matters because regularity is a purely local,
pointwise condition, while the equivalence of the full \emph{solution spaces}
of the Lagrangian and Hamiltonian theories is a global statement about
$\mathbb{F}L$ \cite{abrahammarsden}.

A Lagrangian can be regular everywhere and still fail to be hyperregular, in
either of two ways. First, $\mathbb{F}L$ can fail to be \emph{surjective}: if
$L$ grows sub-quadratically in $\dot q$ at large velocity (rather than the
quadratic-plus-lower-order growth of the usual kinetic-minus-potential form),
the momenta $p=\partial L/\partial\dot q$ can be bounded even as $\dot
q\to\infty$, so that some momentum values in $T^*Q$ are simply never reached,
and the corresponding Hamiltonian trajectories have no Lagrangian counterpart.
Second, $\mathbb{F}L$ can fail to be \emph{injective} globally even while being
a local diffeomorphism everywhere---the fiber derivative can ``fold back'' on
itself far from any particular point, so that two distinct velocities at the
same base point map to the same momentum. In either case, one is forced either
to restrict attention to a proper open subset of phase space on which
$\mathbb{F}L$ does behave well, or to accept that the Lagrangian and
Hamiltonian pictures describe genuinely different (if overlapping) collections
of physically possible histories. This is a comparatively mild failure mode
next to subsections~\ref{singular}--\ref{helmholtz}, but it shows that even the
textbook regularity condition, usually checked (if at all) by verifying $\det
W_{ij}\neq 0$ at a single representative point, is not by itself sufficient for
the global equivalence the folklore claims.

\subsubsection{The Hamilton--Jacobi correspondence: caustics}
\label{caustics}

Jacobi's theorem (Proposition~\ref{prop:jacobi-thm}) is, like
Proposition~\ref{prop:regular}, a genuine equivalence result---but, like
Proposition~\ref{prop:regular} without hyperregularity, it is only a
\emph{local} one, and for exactly the same underlying reason as
subsection~\ref{jacobi}: the nondegeneracy hypothesis on the mixed Hessian
$\partial^2S/\partial q\partial\alpha$ is not guaranteed to persist away from
the point where it is first checked.

Construct Hamilton's principal function concretely, as in
subsection~\ref{hjtheory}: fix an initial event $(q_0,t_0)$, and for each
nearby $(q,t)$ let $S(q,t)$ be the action along the trajectory that solves
Hamilton's equations with $q(t_0)=q_0$ and $q(t)=q$. For $S$ to be a
well-defined, single-valued, smooth function of $(q,t)$, this connecting
trajectory must exist and be \emph{unique}. Existence is guaranteed for short
enough time by ordinary differential equation theory. Uniqueness is exactly
where the conjugate points of subsection~\ref{jacobi} re-enter: consider the
family of trajectories leaving $q_0$ at $t_0$ with varying initial velocity,
and ask where two members of this family, with infinitesimally different
initial velocity, next intersect. This is precisely the kinetic-focus
construction of subsection~\ref{jacobi}, and it recurs at the same time $t_c$
computed there. Beyond $t_c$, distinct nearby trajectories from $(q_0,t_0)$ can
reconverge near the same later point, so more than one candidate value competes
to be $S(q,t)$: the function becomes multivalued, its would-be smooth branches
meeting along an envelope---the \emph{caustic}---where the Jacobian of the map
from initial velocity to endpoint, and with it the mixed Hessian of
Proposition~\ref{prop:jacobi-thm}, degenerates.

\begin{proposition}
\label{prop:caustic}
Let $t_c$ be the first conjugate time to $t_0$ along a trajectory, in the sense
    of subsection~\ref{jacobi}. Hamilton's principal function $S(q,t)$, built
    from trajectories issuing from $(q_0,t_0)$, is a smooth, single-valued
    solution of the Hamilton--Jacobi equation~\eqref{HJ} for $t<t_c$, and
    generically fails to remain single-valued for $t>t_c$, where the family of
    trajectories develops a caustic.
\end{proposition}

Nothing dynamically singular happens to the particle at a caustic: the actual
trajectories---the characteristics of~\eqref{HJ}---pass through it perfectly
smoothly, exactly as they pass through the conjugate point of
subsection~\ref{jacobi}. What breaks down is the specific representational
device of Hamilton--Jacobi theory: a single, smooth, globally valid generating
function from which the entire flow can be read off via $p=\partial S/\partial
q$. This is the same phenomenon, and the same underlying mathematics, that
produces caustics in geometric optics---where Hamilton first developed the
parallel theory---and it is the classical root of the phase corrections (the
Maslov index) that enter semiclassical, WKB-type quantization once a trajectory
has passed through one or more foci \cite{gray2009}.

The upshot mirrors subsection~\ref{jacobi} exactly. ``The Hamiltonian and
Hamilton--Jacobi formulations are equivalent'' is true, without qualification,
only up to the first conjugate point; beyond it, recovering the dynamics from
Hamilton--Jacobi theory requires patching together multiple local branches of
$S$, or abandoning the single smooth generating function in favor of directly
integrating~\eqref{Ham}. As in subsection~\ref{jacobi}, this has nothing to do
with singular Lagrangians, nonholonomic constraints, or a failure of
hyperregularity: the same one-dimensional particle in a smooth potential well
that illustrated subsection~\ref{jacobi} illustrates this failure too, because
a conjugate point in the variational sense and a point on a caustic in the
Hamilton--Jacobi sense are the same event, described in two different
formulations of classical mechanics.

\subsubsection{Non-uniqueness: the correspondence is not one-to-one in the other direction}
\label{nonunique}

Even restricting entirely to the well-behaved case---holonomic constraints, $L$
regular and hyperregular---the map from ``the physics'' to ``a Lagrangian (or
Hamiltonian) representing it'' is many-to-one rather than one-to-one, which
complicates any claim that a given Lagrangian and a given Hamiltonian
formulation are simply \emph{the} representations of one theory. Adding a total
time derivative, $L \mapsto L + \frac{d}{dt}F(q,t)$ for arbitrary $F$, leaves
the Euler--Lagrange equations, and hence all physical predictions, unchanged;
more drastically, so-called $s$-equivalent Lagrangians can differ by more than
a total derivative and still generate the same force law, again governed by the
multiplier version of the Helmholtz analysis of subsection~\ref{helmholtz}. On
the Hamiltonian side, canonical transformations relate infinitely many
distinct-looking triples $(q,p,H)$ that all describe the same underlying
dynamics. None of this contradicts Proposition~\ref{prop:regular}---the
\emph{solution sets} still correspond bijectively---but it means the intuitive
picture of ``one Newtonian system, one Lagrangian, one Hamiltonian''
understates how much freedom is built into the correspondence in both
directions.

\subsection{Discussion: What Should ``Equivalence'' Mean?}
\label{discussion}

subsection~\ref{failures} shows that ``classical mechanics has several
equivalent formulations'' is at best true with qualifications attached, and at
worst equivocates between several distinct things one might mean by
\emph{equivalent}. It is worth separating out at least four senses.

\textbf{Empirical equivalence.} Two formulations are empirically equivalent for
a given system if they issue the same predictions for its observable motion.
This is precisely what fails between nonholonomic and vakonomic mechanics
(subsection~\ref{nonholonomic}): both are internally coherent variational
theories, but experiment sides with one and not the other
\cite{lewismurray1995}.

\textbf{Mathematical (extensional) equivalence.} Two formulations are
extensionally equivalent for a class of systems if there is a bijection between
their solution sets, as in Proposition~\ref{prop:regular}. This is the sense of
equivalence textbooks generally have in mind, and it holds precisely under the
regularity/hyperregularity hypotheses identified in
subsection~\ref{hypotheses}---hypotheses violated by singular Lagrangians
(subsection~\ref{singular}) and inapplicable outright when no Lagrangian exists
at all (subsection~\ref{helmholtz}).

\textbf{Variational equivalence.} A still more demanding sense of equivalence
asks whether the specific mathematical objects that two formulations use to
represent a trajectory---not merely the trajectory itself---continue to exist
and behave as advertised. subsections~\ref{jacobi}--\ref{bvp}
and~\ref{caustics} show this can fail even where Proposition~\ref{prop:regular}
and Proposition~\ref{prop:jacobi-thm} hold without qualification: the true
trajectory satisfies Hamilton's principle, and a complete integral of the
Hamilton--Jacobi equation genuinely generates it, but the further claims that
the trajectory \emph{minimizes} the action, that the two-endpoint problem
Hamilton's principle is actually posed as has exactly one solution, and that
the generating function remains \emph{single-valued}, all fail beyond a
conjugate point that is intrinsic to the potential and can occur arbitrarily
soon. Because all three failures reduce to zeros of the same Jacobi accessory
equation, they are really one failure viewed three ways---once from the
variational side, once from the boundary-value side, once from the
Hamilton--Jacobi side---and all three occur for the least exotic system
imaginable, a single particle in a smooth potential well. This sense of
inequivalence sits below the structural disagreement discussed next: it is not
a dispute about how much geometry a formulation posits, but a demonstration
that some of the load-bearing objects a formulation posits---a minimizing path,
a unique connecting trajectory, a global smooth generating function---simply do
not exist beyond a computable, and often short, horizon.

\textbf{Structural equivalence.} Two formulations are structurally equivalent
if they posit the same amount and kind of mathematical structure on the space
of physical possibilities, even where their solution sets coincide exactly.
This is the level at which a genuine, ongoing disagreement exists in the
philosophy of physics literature. \cite{north2009} argues that Lagrangian and
Hamiltonian mechanics are \emph{not} structurally equivalent even in the fully
regular case: for a ``simple'' mechanical Lagrangian of kinetic-minus-potential
form, the natural state space $TQ$ inherits a Riemannian metric from the
kinetic energy term, whereas the natural Hamiltonian state space $T^*Q$ carries
only a symplectic (volume-like) structure---strictly less structure, on North's
accounting, than a metric. She takes this asymmetry to indicate a genuine
inequivalence between the theories, and argues elsewhere that standard
Newtonian mechanics posits even more structure still, yielding a three-tiered
hierarchy running Newtonian~$>$~Lagrangian~$>$~Hamiltonian \cite{north2022}.
\cite{barrett2015} responds that neither of North's two arguments for the
structural asymmetry succeeds as stated, since compensating structure can also
be identified on the Hamiltonian side; \cite{curiel2014} argues independently,
from the standpoint of constrained and relativistic systems, that Lagrangian
mechanics better represents physical structure than Hamiltonian mechanics does.
\cite{barrett2019} shows that under one precise formal criterion---categorical
equivalence of the two theories construed as structured categories of
models---Lagrangian and Hamiltonian mechanics \emph{are} equivalent, but only
for hyperregular systems, which directly reconnects the philosophical debate to
the hyperregularity condition of subsection~\ref{global}: even the
``equivalent'' verdict is conditional on restricting to the well-behaved class
of systems, and the restriction is not idle, since it excludes precisely the
gauge theories discussed in subsection~\ref{singular}.

The point of separating these four senses is not to declare a winner. It is
that ``are the formulations of classical mechanics equivalent?'' is not a
single question with a single yes-or-no answer. For a specific system and a
specific pair of formulations, one first has to ask which kind of equivalence
is at stake, because the mathematics of subsection~\ref{failures} answers the
empirical and extensional versions of the question rather differently from
system to system, and reasonable, technically informed people disagree about
the structural version even after the mathematics is settled.

\subsection{Conclusion}
\label{conclusion}

The claim that Newtonian, Lagrangian, and Hamiltonian mechanics are
interchangeable formulations of one theory is a good first approximation and a
useful pedagogical simplification, but it is not exactly true. It holds cleanly
only for systems with holonomic constraints (or none), described by a
hyperregular Lagrangian, and even then only in the mathematical/extensional
sense of subsection~\ref{discussion}---a sense some philosophers of physics
contest even in that favorable case. Outside that domain, the failures are not
uniform in character. Some are comparatively benign bookkeeping issues: the
non-uniqueness of subsection~\ref{nonunique}, or the global subtleties of
subsection~\ref{global}, can usually be repaired by working on a suitable open
set or keeping track of gauge equivalence classes. Others are more serious:
singular Lagrangians (subsection~\ref{singular}) are not an edge case but the
generic situation for gauge field theories, meaning that for most of
fundamental physics the naive Legendre-transform equivalence simply does not
apply as stated, and the full Dirac--Bergmann apparatus is doing real work. A
third kind of failure, identified in subsections~\ref{jacobi}, \ref{bvp},
and~\ref{caustics}, needs no exotic system at all: for an ordinary particle in
a smooth potential well, the true trajectory satisfies Hamilton's principle and
generates a Hamilton--Jacobi solution exactly as advertised, yet the further
claims that the action is literally \emph{minimized}, that the two-point
boundary-value problem Hamilton's principle poses has a unique solution, and
that the generating function is globally single-valued, all fail past a
conjugate point set by the curvature of the potential---a point that can be
pushed arbitrarily close to the start by choosing a sufficiently steep
potential well, and beyond which the boundary-value problem can have no
solution at all or infinitely many, a multiplicity that resurfaces quantum
mechanically as the sum over classical paths and the associated Maslov index.
And one failure is not a bookkeeping issue at all: the split between
nonholonomic and vakonomic mechanics (subsection~\ref{nonholonomic}) is a case
where two internally consistent variational formulations of the same
constrained system make different predictions, decided by experiment rather
than by mathematical convention. That a subject as old and as thoroughly
axiomatized as classical mechanics still contains a live fork of this kind is a
useful reminder that ``equivalent formulations of the same physics'' is a claim
to be checked, case by case, rather than assumed.
